%% notes-tex/033-post-query-analysis/sections/1-question.tex

\section{The question\secstatus{todo}}
\label{sec:question}

The planning documents for inhibitor-tolerance modeling,
\texttt{notes-tex/031-inhibitor-tolerance-data} and
\texttt{notes-tex/031-unified-representation}, measured everything on the dev-tree
loaders, one row per served record, and most statistics on the (gene, compound) pair. A
model does not train on those rows. It trains on the store the query built, build 001 of
experiment 033, whose row is a strain in one environment. Whether the plan's numbers
describe that store is an empirical question, and every answer here is read from the store
itself.

Every script named below lives in \file{experiments/033-env-chemgen-pooled/scripts/}.
\file{flatten_cells.py} reads each entry of the processed store once and writes its
genotype, environment and every folded measurement as one table row (slurm 2988).
\file{store_against_plan.py} re-measures each plan claim on that table and reads the
plan's own values from the 031 result files at commit \texttt{bf1884c46}.
\file{post_query_figures.py} and \file{notes_tex_post_query_tables.py} draw the figures
and render the tables.

\notesec{Vocabulary}
A \term{cell} is one entry of the store: a genotype, its ploidy, and one environment. An
\term{environment} is the content address the graph adapter gives a condition: medium
composition, temperature, the dosed compounds at their concentrations, aerobicity and
duration. The \term{functional dose} of a perturbed gene is the copies present over the
copies in the reference, zero for a deletion and one half for a heterozygous deletion in a
diploid. The \term{label} of a cell is the one number the store's label table holds for it.
A \term{measurement} is one served record, and a cell folds every measurement that shares
its genotype and environment. The \term{orientation} of a source is the factor, plus or
minus one, that makes a negative response a fitness defect. A \term{pair} is one (gene,
compound) combination with the response averaged over every dose, screen and ploidy arm
that measured it, which is the unit the plan's joinability statistics use. The \term{plan}
is the two 031 documents together. The \term{reliability} of a served value is the Pearson
correlation over genes between two screens of the same environment, and its square root is
the \term{ceiling} on the correlation any model can reach with the noise-free response. The
\term{store} is build 001 of the pooled chemogenomic dataset,
\file{/db/experiments/033-env-chemgen-pooled-001-pooled-build}.

\begin{keybox}
\textbf{Measured.} Every axis of the store per source (Section~\ref{sec:units}); the
ploidy arms and what pairs them (Section~\ref{sec:genotype}); the environment vocabulary,
the dose and the compound vectors (Section~\ref{sec:environment}); which measurement the
label table holds and how well repeated screens agree (Section~\ref{sec:label}); the pooled
target after orienting and standardizing (Section~\ref{sec:target}).

\textbf{Not measured.} No model has been trained on the store. Every statement about what
a loss, a token or an operator would do with these numbers is labeled as a hypothesis, and
Section~\ref{sec:next} lists the decisions the measurements leave open.
\end{keybox}
